\section{The volume formulation}
\label{sec:vol}

Gila assembles every entry of the discretized vacuum Green operator from thirty-six
face-pair integrals of the scalar kernel. That assembly is a signed sum of nearly equal
numbers, and it is the reason the far field is only good to $10^{-10}$ today
(Section~\ref{sec:vol:pin} quantifies it). This section undoes it: it shows that the
signed sum of the thirty-six face pairs \emph{is} a single smooth integral over the
difference box of the pair, derives every constant and every sign of that identity from
the code rather than asserting them, checks the identity in exact rational arithmetic on
the polynomial kernels $|\mathbf{x}-\mathbf{y}|^{2j}$, and pins it numerically against a
$220$-bit reference. Every expansion in the rest of this document expands that one
integral.

\subsection{Conventions, and what \code{srfMat} holds}
\label{sec:vol:conv}

Lengths are in wavelengths, the frequency is the single complex scalar $f$ (the field
\code{frqPhz}), $k = 2\pi f$, and Gila's scalar kernel is
\begin{equation}
g(r) \;=\; \frac{e^{2\pi\ii f r}}{4\pi f^{2} r} \;=\; \frac{e^{\ii k r}}{4\pi f^{2} r},
\qquad r = |\mathbf{x}-\mathbf{y}| .
\label{eq:vol:ker}
\end{equation}
A cell has edge lengths $\mathbf{s} = (s_1,s_2,s_3)$, exact \code{Rational}s in
\code{GlaVol.scl}; the source cell $V_{\mathrm{s}}$ is centred at the origin, the target
cell $V_{\mathrm{t}}$ at $\mathbf{R} = (n_1 s_1, n_2 s_2, n_3 s_3)$ with integer
$\mathbf{n}$, and $V_{\mathrm{t}} = s_1 s_2 s_3$ denotes the cell volume as well as the
target cell. The six faces are ordered \code{yzL yzU xzL xzU xyL xyU}; face $F$ has axis
$c(F) \in \{1,2,3\}$ and outward normal $\sigma_F \mathbf{e}_{c(F)}$ with
$\sigma_F = -1$ for the lower (L) and $+1$ for the upper (U) face. The face-pair index is
$\mathrm{fp} = 6(F-1) + F'$ with $F$ the \emph{target} face and $F'$ the \emph{source}
face (\code{facPar} writes \code{fPairs[1,k] = i} for the target and
\code{fPairs[2,k] = j} for the source).

\begin{lemma}[what the code stores]
\label{lem:vol:scl}
Let $I_{FF'} = \int_F\int_{F'} g\,\dd S'\,\dd S$ be the raw face-pair integral. Then
\code{srfMat[fp]} $= I_{FF'}/V_{\mathrm{t}}$.
\end{lemma}

\begin{proof}
\code{egoSrfFxd!} evaluates the integrand on the unit hypercube $[0,1]^4$ with the
bilinear maps of the two rectangles built by \code{cubFac} and multiplies the result by
\code{srfScales[fp]}. The bilinear map of a rectangle $F$ onto $[0,1]^2$ has constant
Jacobian $A_F$, the area of $F$, so the quadrature approximates
$\int_{[0,1]^4} = I_{FF'}/(A_F A_{F'})$. \code{srfScl} sets
\code{srfScl[(trgFId-1)*6+srcFId]} to $A_{F'}/s_{c(F)}$: the source face area over the
target edge normal to the target face. Hence
$\code{srfMat[fp]} = I_{FF'}/(A_F A_{F'}) \cdot A_{F'}/s_{c(F)} = I_{FF'}/(A_F\,s_{c(F)})$,
and for a cuboid $A_F\, s_{c(F)} = V_{\mathrm{t}}$ for every face $F$.
\end{proof}

\subsection{The divergence theorem, read backwards}
\label{sec:vol:div}

Equation~(7) of the companion document \emph{The Integrals of Gila} turns the
volume--volume integral of the strongly singular dyadic into surface--surface integrals of
the weakly singular kernel. Read in the other direction it recovers a volume integral,
and that is the direction the far field needs.

\begin{lemma}[the divergence identity in both directions]
\label{lem:vol:div}
Let $F$ be $C^2$ on a neighbourhood of the (closed) difference set
$V_{\mathrm{t}} - V_{\mathrm{s}}$. Then for $a \neq b$
\begin{equation}
\int_{V_{\mathrm{t}}}\!\int_{V_{\mathrm{s}}}
(\partial_a \partial_b F)(\mathbf{x}-\mathbf{y})\,\dd\mathbf{y}\,\dd\mathbf{x}
\;=\; -\sum_{\substack{F:\,c(F)=a\\ F':\,c(F')=b}} \sigma_F\,\sigma_{F'}\,
\int_F\!\int_{F'} F\,\dd S'\,\dd S ,
\label{eq:vol:divoff}
\end{equation}
and for the diagonal, using
$(\partial_a\partial_a - \nabla^2)F = -\sum_{b \neq a}\partial_b\partial_b F$,
\begin{equation}
\int_{V_{\mathrm{t}}}\!\int_{V_{\mathrm{s}}}
\bigl[(\partial_a \partial_a - \nabla^2) F\bigr](\mathbf{x}-\mathbf{y})\,
\dd\mathbf{y}\,\dd\mathbf{x}
\;=\; +\sum_{b \neq a}\ \sum_{\substack{F,F':\\ c(F)=c(F')=b}} \sigma_F\,\sigma_{F'}\,
\int_F\!\int_{F'} F\,\dd S'\,\dd S .
\label{eq:vol:divdia}
\end{equation}
\end{lemma}

\begin{proof}
$F$ depends on $\mathbf{x}$ and $\mathbf{y}$ only through $\mathbf{u} =
\mathbf{x}-\mathbf{y}$, so $\partial_{u_a} = \partial^{\mathbf{x}}_a =
-\partial^{\mathbf{y}}_a$ acting on $F(\mathbf{x}-\mathbf{y})$, whence
$(\partial_a\partial_b F)(\mathbf{x}-\mathbf{y}) = -\partial^{\mathbf{x}}_a
\partial^{\mathbf{y}}_b F(\mathbf{x}-\mathbf{y})$. Each derivative is moved onto its own
domain by the divergence theorem, $\int_{V}\partial_a\Phi\,\dd\mathbf{x} =
\oint_{\partial V} n_a \Phi\,\dd S$, giving
$-\oint_{\partial V_{\mathrm{t}}}\oint_{\partial V_{\mathrm{s}}}
n_a(\mathbf{x})\,n'_b(\mathbf{y})\,F\,\dd S'\,\dd S$. On a cuboid $n_a$ vanishes on the
four faces whose axis is not $a$ and equals $\sigma_F$ on the two whose axis is $a$, which
is \eqref{eq:vol:divoff}. The diagonal follows by applying \eqref{eq:vol:divoff} with
$b = a$ to each of the two terms of
$-\sum_{b\neq a}\partial_b\partial_b F$, the two minus signs cancelling.
\end{proof}

\begin{proposition}[\code{srfSum!} is exactly Lemma~\ref{lem:vol:div}]
\label{prop:vol:srfsum}
With $\code{srfMat[fp]} = I_{FF'}/V_{\mathrm{t}}$ of Lemma~\ref{lem:vol:scl}, the nine
signed sums of \code{srfSum!} (\code{glaVacOprMemGen.jl}, lines 734--757) are
\begin{equation}
\bigl(G_{\mathrm{gila}}\bigr)_{ab}(\mathbf{R})
\;=\; \frac{1}{V_{\mathrm{t}}}\int_{V_{\mathrm{t}}}\!\int_{V_{\mathrm{s}}}
\bigl[(\partial_a\partial_b - \delta_{ab}\nabla^2)\,g\bigr](\mathbf{x}-\mathbf{y})\,
\dd\mathbf{y}\,\dd\mathbf{x},
\label{eq:vol:srfsum}
\end{equation}
with the sign $+1$ on the diagonal and on the off-diagonals and with no further factor of
$k^2$, $4\pi$, $f^2$ or $V_{\mathrm{t}}$.
\end{proposition}

\begin{proof}
Read the coefficients off the code and compare with Lemma~\ref{lem:vol:div}. For the
off-diagonal entry $(1,2)$, \code{srfSum!} gives
$-\code{srfMat[3]}+\code{srfMat[4]}+\code{srfMat[9]}-\code{srfMat[10]}$; the four indices
are $\mathrm{fp} = 6(F-1)+F'$ for $(F,F') = (1,3),(1,4),(2,3),(2,4)$, i.e.\ the two
$x$-faces of the target against the two $y$-faces of the source, and their normal signs
are $\sigma_F\sigma_{F'} = (-)(-) = +1$, $(-)( +) = -1$, $(+)(-) = -1$, $(+)(+) = +1$. The
code's coefficients are $-1, +1, +1, -1$, i.e.\ exactly $-\sigma_F\sigma_{F'}$, which is
\eqref{eq:vol:divoff}. For the diagonal entry $(1,1)$, \code{srfSum!} gives
$\code{srfMat[15]}-\code{srfMat[16]}-\code{srfMat[21]}+\code{srfMat[22]}
+\code{srfMat[29]}-\code{srfMat[30]}-\code{srfMat[35]}+\code{srfMat[36]}$; those eight
indices are $(F,F') = (3,3),(3,4),(4,3),(4,4)$ (both faces on the $y$ axis) and
$(5,5),(5,6),(6,5),(6,6)$ (both on the $z$ axis), with
$\sigma_F\sigma_{F'} = +1,-1,-1,+1$ in each block, matching the code's coefficients with
no extra sign, which is \eqref{eq:vol:divdia}. The remaining seven entries are the same
two patterns with the axes permuted.
\end{proof}

\subsection{Difference variables and the triangle weight}
\label{sec:vol:diff}

\begin{lemma}[the difference box]
\label{lem:vol:box}
For any integrable $\Phi$,
\begin{equation}
\int_{V_{\mathrm{t}}}\!\int_{V_{\mathrm{s}}} \Phi(\mathbf{x}-\mathbf{y})\,
\dd\mathbf{y}\,\dd\mathbf{x}
\;=\; \int_{D} w(\boldsymbol{\delta})\,\Phi(\mathbf{R}+\boldsymbol{\delta})\,
\dd\boldsymbol{\delta},
\qquad
D = \prod_{i=1}^{3}[-s_i, s_i],
\quad
w(\boldsymbol{\delta}) = \prod_{i=1}^{3}\bigl(s_i - |\delta_i|\bigr).
\label{eq:vol:diffbox}
\end{equation}
\end{lemma}

\begin{proof}
The integrand depends only on $\mathbf{u} = \mathbf{x}-\mathbf{y}$ and both cells are
products of intervals, so the identity factors over the axes. On axis $i$ the target
interval is $R_i + [-s_i/2, s_i/2]$ and the source interval is $[-s_i/2, s_i/2]$; the
push-forward of the product Lebesgue measure under $(x_i,y_i)\mapsto x_i - y_i$ is the
convolution of the two indicators, a triangle of base $2s_i$ and apex height $s_i$
centred at $R_i$. Writing $u_i = R_i + \delta_i$ gives the density $s_i - |\delta_i|$ on
$[-s_i,s_i]$.
\end{proof}

The moments of that weight are the only geometry the far field ever needs:
\begin{equation}
\mu_n(s) \;=\; \int_{-s}^{s}(s-|t|)\,t^n\,\dd t
\;=\; \begin{cases} \dfrac{2\,s^{\,n+2}}{(n+1)(n+2)}, & n \text{ even},\\[4pt]
0, & n \text{ odd},\end{cases}
\qquad
\int_D w\,\boldsymbol{\delta}^{\boldsymbol{\alpha}}\,\dd\boldsymbol{\delta}
= \prod_{i=1}^{3}\mu_{\alpha_i}(s_i),
\label{eq:vol:mu}
\end{equation}
by $\mu_n = 2\int_0^s (s-t)t^n\,\dd t$ for even $n$ and oddness of the integrand
otherwise. Note $\mu_0(s) = s^2$, so $\int_D w = V_{\mathrm{t}}^{\,2}$ and the cell
average $(1/V_{\mathrm{t}})\int_D w = V_{\mathrm{t}}$.

Combining Proposition~\ref{prop:vol:srfsum}, Lemma~\ref{lem:vol:box} and the Helmholtz
equation $\nabla^2 g = -k^2 g$, valid away from $r = 0$ and therefore on all of
$\mathbf{R}+D$ as soon as the two cells do not touch, gives the object of this document.

\begin{theorem}[the volume form of the far field]
\label{thm:vol:main}
Let the two cells be separated, i.e.\ $\mathbf{0} \notin \mathbf{R}+D$. Then
\begin{equation}
\boxed{\;
\begin{aligned}
\bigl(G_{\mathrm{gila}}\bigr)_{ab}(\mathbf{R})
&\;=\; \frac{1}{V_{\mathrm{t}}}\int_{D} w(\boldsymbol{\delta})\,
\bigl[(\partial_a\partial_b + \delta_{ab}k^2)\,g\bigr]
(\mathbf{R}+\boldsymbol{\delta})\,\dd\boldsymbol{\delta}\\
&\;=\; \frac{1}{V_{\mathrm{t}}}\int_{D} w(\boldsymbol{\delta})\,
\bigl[(\partial_a\partial_b - \delta_{ab}\nabla^2)\,g\bigr]
(\mathbf{R}+\boldsymbol{\delta})\,\dd\boldsymbol{\delta} .
\end{aligned}\;}
\label{eq:vol:main}
\end{equation}
The integrand is analytic on a neighbourhood of $D$; explicitly, with
$\mathbf{z} = \mathbf{R}+\boldsymbol{\delta}$, $r = |\mathbf{z}|$,
\begin{equation}
g' = \Bigl(\ii k - \frac{1}{r}\Bigr) g, \qquad
g'' = \Bigl(\ii k - \frac{1}{r}\Bigr) g' + \frac{g}{r^2}, \qquad
\partial_a\partial_b\, g = (g'' - g'/r)\,\frac{z_a z_b}{r^2} + \delta_{ab}\,\frac{g'}{r},
\label{eq:vol:d2}
\end{equation}
and $\nabla^2 g = g'' + 2g'/r = -k^2 g$.
\end{theorem}

The point of \eqref{eq:vol:main} is that the right-hand side carries no cancelling signed
sum: $w \ge 0$ on $D$, and the thirty-six face-pair integrals that Gila subtracts have
been replaced by one integral of a smooth function against a non-negative weight. What
the far field must do is evaluate that integral without quadrature.

\subsection{Relation to the physical dyadic}
\label{sec:vol:phys}

The free-space dyadic and the form Gila uses are
\begin{equation}
\mathbf{G}(\mathbf{r}) = \Bigl(\mathbf{I} + \frac{\nabla\nabla}{k^2}\Bigr)\Gs(r)
= \frac{1}{k^2}\bigl(\nabla\nabla - \mathbf{I}\nabla^2\bigr)\Gs(r)
- \frac{\mathbf{I}}{k^2}\delta(\mathbf{r}),
\qquad \Gs(r) = \frac{e^{\ii k r}}{4\pi r},
\label{eq:vol:phys}
\end{equation}
and the entry of the discretized operator is the double cell average
$(1/V_{\mathrm{t}})\int_{V_{\mathrm{t}}}\int_{V_{\mathrm{s}}} G_{ab}$. Gila's operator is
$4\pi^2$ times that average in units where the wavelength is one; the fixed constant is
visible in \code{genEgoCrcSlf!}, where the identity term $-\delta_{ab}\delta_{ts}/k^2$ is
implemented as \code{egoToe[i,i,1,1,1] -= 1/frqPhz\^{}2}, i.e.\ as $-1/f^2 =
-4\pi^2/k^2$. The two constants of \eqref{eq:vol:phys} and of Gila's normalization
combine once and for all:
\begin{equation}
\frac{4\pi^2}{k^2} = \frac{4\pi^2}{4\pi^2 f^2} = \frac{1}{f^2},
\qquad\text{so}\qquad
\frac{4\pi^2}{k^2}\,\Gs(r) = \frac{e^{\ii k r}}{4\pi f^2 r} = g(r).
\label{eq:vol:const}
\end{equation}
\emph{The $1/k^2$ of the dyadic and the $4\pi^2$ of Gila's convention are exactly the
$1/(4\pi f^2)$ already inside $g$}, which is why no $1/k^2$ and no $4\pi^2$ survive
anywhere in \eqref{eq:vol:main}: they are not dropped, they are absorbed. Two
consequences are worth recording because they are cheap end-to-end tests. First, every
integral of $1/r$ is proportional to $f^{-2}$ and nothing else, so $g$ and hence
$G_{\mathrm{gila}}$ is homogeneous, $G_{\mathrm{gila}}(\alpha\mathbf{s},
\alpha\mathbf{n}, f/\alpha) = \alpha^{2}\,G_{\mathrm{gila}}(\mathbf{s},\mathbf{n},f)$.
Second, the far entries are small: at $\lambda/32$, $f=1$, the largest entry of
$G_{\mathrm{gila}}$ on an axis offset is $2.08\times10^{-2}$ at two cells,
$5.77\times10^{-4}$ at eight, $9.44\times10^{-5}$ at thirty-two and
$2.39\times10^{-5}$ at $128$ cells, against a self term of order one (the identity
$-1/f^2$). Accuracy ``relative to the largest entry of the tensor'' is therefore four
orders looser than per-entry accuracy for a far offset, and per-entry relative accuracy in
\code{Float64} is the honest target throughout this document.

\subsection{The identity in exact rational arithmetic}
\label{sec:vol:exact}

Theorem~\ref{thm:vol:main} contains two claims that can be separated and tested exactly:
the divergence identity of Lemma~\ref{lem:vol:div} with the \code{srfSum!} signs, and the
difference-box reduction of Lemma~\ref{lem:vol:box}. On a polynomial kernel both sides are
rational numbers in the edge lengths and can be compared with no arithmetic error at all.
Take $F(\mathbf{u}) = |\mathbf{u}|^{2j}$. Then
$(\partial_a\partial_a - \nabla^2)|\mathbf{u}|^{2} = -4$ identically, and more generally
every side of \eqref{eq:vol:divoff}--\eqref{eq:vol:divdia} is a finite rational
combination of the $\mu_n(s_i)$ and of the corresponding face moments.

The check was run for $\mathbf{n} \in \{(2,0,0),(3,1,0),(2,2,2)\}$ and $j = 1,2,3$ at
$\mathbf{s} = (1/32)^3$ in \code{Rational} of \code{BigInt} arithmetic, and both identities hold
\emph{exactly} --- not to a tolerance. The volume moments $\int_D w\,
|\mathbf{R}+\boldsymbol{\delta}|^{2j}\,\dd\boldsymbol{\delta}$, computed by two
independent exact routes (multinomial expansion against \eqref{eq:vol:mu}, and the
face-pair route), agree exactly and are

\begin{center}\small
\begin{tabular}{@{}lccc@{}}
\toprule
$\mathbf{n}$ & $j = 1$ & $j = 2$ & $j = 3$\\
\midrule
$(2,0,0)$ & $\dfrac{9}{2199023255552}$ & $\dfrac{691}{33776997205278720}$
          & $\dfrac{82279}{726340547902313594880}$\\[6pt]
$(3,1,0)$ & $\dfrac{21}{2199023255552}$ & $\dfrac{3511}{33776997205278720}$
          & $\dfrac{866881}{726340547902313594880}$\\[6pt]
$(2,2,2)$ & $\dfrac{25}{2199023255552}$ & $\dfrac{4931}{33776997205278720}$
          & $\dfrac{285155}{145268109580462718976}$\\
\bottomrule
\end{tabular}
\end{center}

\noindent
The smallest case can be checked by hand and is worth writing out, because it fixes the
sign of \eqref{eq:vol:divdia} with no computer at all: for $j = 1$,
$(\partial_1\partial_1 - \nabla^2)|\mathbf{u}|^2 = 2 - 6 = -4$, so the left-hand side of
\eqref{eq:vol:divdia} is $-4\int_D w = -4\prod_i \mu_0(s_i) = -4 s^6 = -4/32^6 =
-1/268435456$ at $\mathbf{s} = (1/32)^3$, and the signed sum of the corresponding
face-pair moments is that same rational number. The even moments of
\code{notes/moments/moments.jl} agree with the exact polynomial integrator to $154$
digits at $512$ bits, which is that library's own floor.

\subsection{The identity pinned numerically, and the sign}
\label{sec:vol:pin}

The exact check of Section~\ref{sec:vol:exact} tests the geometry but not the kernel. The
kernel test compares two $220$-bit evaluations of the same tensor that share no code: the
\emph{face-pair reference} (the thirty-six integrals $I_{FF'}$ of $g$ from \code{pairKer}
of \code{notes/gen/ref/mom.jl} at \code{ordN} $=44$, divided by $V_{\mathrm{t}}$ and
assembled with the \code{srfSum!} signs) and the \emph{volume reference} (a graded
\code{BigFloat} tensor Gauss--Legendre rule of order $32$ applied to the right-hand side
of \eqref{eq:vol:main}, with every axis of $D$ split at $0$ where $w$ has its kink, and
panels graded geometrically towards the corner nearest the singularity at
$\boldsymbol{\delta} = -\mathbf{R}$; nodes generated by Newton iteration on the Legendre
recursion in \code{BigFloat}, so nothing in the rule is \code{Float64}). All four sign
conventions were tested at every offset: $e_{++}$ is $(+\text{diag},+\text{off})$,
$e_{+-}$ is $(+\text{diag},-\text{off})$, and so on, each entry being
$\max_{ab}|G_{\text{conv}} - G_{\text{ref}}|/\max_{ab}|G_{\text{ref}}|$;
$\mathrm{ent}_{++}$ is the worst per-entry relative error of the surviving convention.
Shapes: $c_{32} = (1/32)^3$, $c_8 = (1/8)^3$, $c_4 = (1/4)^3$, $\mathrm{sl} =
(1/32,1/32,1/512)$; $f = 1$ (r) or $1+0.1\ii$ (c).

\begin{center}\footnotesize
\begin{tabular}{@{}llrccccc@{}}
\toprule
shape & $f$ & $\mathbf{n}$ & $e_{++}$ & $\mathrm{ent}_{++}$ & $e_{+-}$ & $e_{-+}$ & $e_{--}$\\
\midrule
$c_{32}$ & r & $(2,0,0)$    & 2.43e-58 & 2.81e-58 & 2.43e-58 & 2.0 & 2.0\\
$c_{32}$ & r & $(2,2,0)$    & 9.72e-62 & 1.74e-61 & 2.0 & 1.115 & 2.0\\
$c_{32}$ & r & $(2,2,2)$    & 2.00e-63 & 3.60e-63 & 2.0 & 0.567 & 2.0\\
$c_{32}$ & r & $(3,1,0)$    & 1.09e-59 & 2.31e-59 & 0.943 & 2.0 & 2.0\\
$c_{32}$ & r & $(8,0,0)$    & 1.16e-63 & 1.16e-63 & 1.16e-63 & 2.0 & 2.0\\
$c_{32}$ & r & $(64,0,0)$   & 2.63e-63 & --- & 2.63e-63 & 2.0 & 2.0\\
$c_{32}$ & r & $(64,64,64)$ & 4.14e-63 & 8.26e-63 & 1.003 & 2.0 & 2.0\\
$c_{32}$ & c & $(2,0,0)$    & 2.44e-58 & 2.77e-58 & 2.44e-58 & 2.0 & 2.0\\
$c_{32}$ & c & $(8,8,8)$    & 6.16e-63 & 8.97e-63 & 1.373 & 2.0 & 2.0\\
$c_{8}$  & r & $(2,0,0)$    & 1.49e-58 & 1.49e-58 & 1.49e-58 & 2.0 & 2.0\\
$c_{8}$  & c & $(4,4,4)$    & 6.80e-64 & 1.22e-63 & 1.111 & 2.0 & 2.0\\
$c_{4}$  & r & $(2,0,0)$    & 7.41e-59 & 1.05e-58 & 7.41e-59 & 2.0 & 2.0\\
$c_{4}$  & r & $(2,2,2)$    & 1.53e-64 & 2.90e-64 & 1.054 & 2.0 & 2.0\\
$c_{4}$  & c & $(16,0,0)$   & 3.71e-65 & --- & 3.71e-65 & 2.0 & 2.0\\
sl & r & $(2,0,0)$   & 1.85e-58 & 4.50e-58 & 1.85e-58 & 2.0 & 2.0\\
sl & r & $(0,2,0)$   & 1.85e-58 & 4.50e-58 & 1.85e-58 & 2.0 & 2.0\\
sl & r & $(0,0,2)$   & 1.05e-49 & 1.07e-49 & 1.05e-49 & 2.0 & 2.0\\
sl & c & $(0,0,16)$  & 9.41e-58 & 9.99e-58 & 9.41e-58 & 2.0 & 2.0\\
\bottomrule
\end{tabular}
\end{center}

\noindent
The $(+\text{diag},+\text{off})$ convention of Theorem~\ref{thm:vol:main} agrees with the
face-pair reference to between $3.7\times10^{-65}$ and $1.05\times10^{-49}$ at all
eighteen offsets, four shapes and both frequencies --- nineteen to thirty-four orders
better than the $10^{-30}$ that was asked for --- and every other convention is wrong by
$0.57$ to $2.0$. \emph{The sign is $+1$ on the diagonal and on the off-diagonals, with no
extra factor.} The rows where $e_{+-} = e_{++}$ are the offsets whose off-diagonal
entries vanish by symmetry; the off-diagonal sign is pinned by the seven rows where
$e_{+-}$ lies between $0.943$ and $2.0$. The two suppressed $\mathrm{ent}_{++}$ values
are an artefact of the per-entry metric at offsets whose off-diagonals vanish
identically, where \code{pairKer} leaves them at its own roundoff
($7.2\times10^{-71}$ and $2.3\times10^{-65}$ absolute) and the quotient is $O(1)$. The
single worst row, slender $(0,0,2)$ at $1.05\times10^{-49}$, is the volume rule's own
truncation and not a disagreement: there the difference box is sixteen times wider than
it is deep and the singularity sits $1/512$ outside one of its faces, so the graded rule
at order $32$ is itself converged only to $1.8\times10^{-24}$, and $250$~s are needed per
evaluation. Everything else sits at or near the $220$-bit floor of $1.3\times10^{-64}$.

Two further independent confirmations were obtained, deliberately without shared code.
A second \code{BigFloat} Gauss--Legendre volume rule, written from scratch, reproduces the
face-pair reference to $10^{-63}$--$10^{-65}$ on the cubes and $1.4\times10^{-62}$ on the
slender cell. And an independent \code{NIntegrate} evaluation of the right-hand side of
\eqref{eq:vol:main} in Mathematica at $\mathbf{n} = (5,1,0)$, $\mathbf{s} = (1/32)^3$,
$f = 1$, working precision $30$--$40$, the box split at $0$: entry $(1,2)$ agrees with the
$220$-bit face-pair reference to $3.4\times10^{-14}$, which is Mathematica's own error
estimate, and entry $(1,1)$, recomputed with a fixed $30$-point Gauss--Kronrod rule per
sub-box after the global adaptive strategy stalled, reproduces
$0.001570738221266370299274 + 0.000360155079207178606379\,\ii$ in every printed digit.
Sign $+$, no extra factor, normalization confirmed by three routes.

\subsection{Why the volume form and not the face pairs}
\label{sec:vol:amp}

The assembly of Proposition~\ref{prop:vol:srfsum} subtracts nearly equal numbers, and the
amplification is measurable. Define
$\Lambda_{ab} = \bigl(\sum_{\mathrm{fp}} |\code{srfMat}_{\mathrm{fp}}|\bigr) /
|G_{ab}|$ over the face pairs summed into entry $(a,b)$. Against the $220$-bit reference,
Gila's present fixed Gauss--Legendre rule achieves a per-face-pair relative error of
$10^{-15}$ to $9\times10^{-14}$ at $\lambda/32$, but $\Lambda$ runs from $36$ to $1300$
(it grows as $n^2$ on an axis offset and is flat at $312$ on the body diagonal), so the
assembled tensor error is $5.0\times10^{-14}$ to $8.7\times10^{-13}$. At $\lambda/8$ the
amplification is $14$--$326$ and the tensor error $5\times10^{-15}$ to
$3\times10^{-13}$; at $\lambda/4$, $5$--$163$ and $1.6\times10^{-15}$ to
$2.8\times10^{-13}$. The amplification is a \emph{fine-cell} effect, growing like
$(|\mathbf{R}|/s)^2$. On the slender cell it is far worse: the long axes carry
$\Lambda = 4.3\times10^{3}$ to $4.0\times10^{4}$, giving $3.2\times10^{-12}$ to
$2\times10^{-11}$ despite per-pair errors of $10^{-15}$ to $8\times10^{-14}$, and along
the short axis the fixed rule is wrong by $O(1)$: relative tensor errors of $0.79$ at
$(0,0,2)$, $3.7\times10^{-2}$ at $(0,0,4)$, $1.8\times10^{-3}$ at $(0,0,8)$,
$3\times10^{-6}$ at $(0,0,16)$, $4.5\times10^{-6}$ at $(0,0,20)$ and
$1.9\times10^{-6}$ at $(0,0,24)$, because \code{quadOrd} keys the rule order on the
separation in \emph{cells} and not on the physical geometry, and a $1/512$ boundary layer
across a $1/32$ face is invisible to it.

Equation~\eqref{eq:vol:main} removes $\Lambda$ from the problem entirely. It is the
formulation used by every family of the registry (Section~\ref{sec:reg}).
